\chapter*{Course Syllabus \& Instructional Design}
\addcontentsline{toc}{chapter}{Course Syllabus \& Instructional Design}

\section*{Course Information}
\begin{itemize}
    \item \textbf{Course Code \& Title:} PHYS2026 Mathematics for Physics
    \item \textbf{Credit Hours:} 3 Credits (3-0-6) [Lecture 3 hrs, Practical 0 hrs, Self-Study 6 hrs/week]
    \item \textbf{Curriculum:} Bachelor of Science in Physics, Core Major Requirement
    \item \textbf{Lead Instructor:} Asst. Prof. Dr. Chewa Thassana (chewa.t@rbru.ac.th)
    \item \textbf{Target Academic Level:} Year 2, Semester 1
\end{itemize}

\section*{Course Description}
Advanced mathematical methods and computational techniques for physical applications. Topics include: curvilinear coordinate systems (Cartesian, cylindrical, spherical) and metric scale factors; vector analysis, gradient, divergence, curl, and conservative fields; Taylor and Maclaurin expansions, asymptotic approximations, and Fourier series; complex numbers, Euler's formula, phasors, and harmonic oscillations; matrix algebra, determinants, eigenvalues, eigenvectors, and diagonalization; multivariable differential calculus, tangent planes, directional derivatives, and Lagrange multipliers; multiple integrals, Jacobians, Gauss's Divergence theorem, and Stokes' theorem; analytical and numerical solutions of first-order and second-order ordinary differential equations with physical modeling in Python.

\section*{Course Learning Outcomes (CLOs) \& Bloom's Taxonomy}
Upon successful completion of this course, students will be able to:
\begin{enumerate}
    \item \textbf{CLO 1 [Cognitive: Understanding / Remembering]:} Articulate foundational mathematical concepts and their direct correspondence to physical laws in classical mechanics, electromagnetism, and wave phenomena.
    \item \textbf{CLO 2 [Cognitive: Applying / Analyzing]:} Formulate physical boundary-value problems in the most symmetrical curvilinear coordinate system and apply vector differential operators accurately.
    \item \textbf{CLO 3 [Cognitive: Analyzing / Evaluating]:} Solve first- and second-order linear differential equations analytically, evaluating asymptotic behaviors, damping regimes, and resonance phenomena.
    \item \textbf{CLO 4 [Practical / Skill: Creating / Computing]:} Implement scientific Python algorithms to compute numerical trajectories, evaluate matrix eigenvalues, and visualize multivariable vector fields.
\end{enumerate}

\section*{Course Curriculum Map (9 Chapters)}
\begin{table}[H]
\centering
\small
\begin{tabularx}{\textwidth}{clp{7.5cm}c}
\toprule
\textbf{Week} & \textbf{Chapter} & \textbf{Topic Overview} & \textbf{Target CLO} \\
\midrule
1--2 & Chapter 1 & Curvilinear Coordinate Systems \& Scale Factors & CLO 1, 2 \\
3--4 & Chapter 2 & Vector Analysis \& Differential Field Operators & CLO 1, 2 \\
5 & Chapter 3 & Infinite Series, Taylor Expansions \& Approximations & CLO 1, 2 \\
6--7 & Chapter 4 & Complex Analysis, Euler's Formula \& Phasors & CLO 1, 3 \\
8--9 & Chapter 5 & Matrices, Eigenvalues \& Coupled Oscillations & CLO 2, 4 \\
10 & Chapter 6 & Multivariable Differential Calculus \& Optimization & CLO 2, 3 \\
11--12 & Chapter 7 & Multiple Integrals \& Integral Theorems (Gauss/Stokes) & CLO 2, 3 \\
13 & Chapter 8 & First-Order Ordinary Differential Equations & CLO 3, 4 \\
14--15 & Chapter 9 & Second-Order ODEs, Damped Oscillations \& Resonance & CLO 3, 4 \\
\bottomrule
\end{tabularx}
\caption{Instructional schedule and curriculum alignment with CLOs.}
\label{tab:syllabus_map}
\end{table}

\section*{Evaluation and Grading Policy}
\begin{itemize}
    \item \textbf{Continuous Assessment (Formative):} 60\%
    \begin{itemize}
        \item Homework Problem Sets (Tiered Conceptual and Analytical): 25\%
        \item Active Learning Python Simulation Mini-Projects: 20\%
        \item Midterm Examination (Chapters 1--5): 15\%
    \end{itemize}
    \item \textbf{Comprehensive Final Examination (Summative):} 40\%
\end{itemize}
